\begin{remark}
	If \( \qo{L}{K} \) is finite, then any morphism of \( K \)-algebra \( L \to L \) are isomorphism.
	\begin{itemize}
		\item It is automatically injective.
		\item We have
		      \[
			      K \subseteq \sigma(L) \subseteq L \implies \sigma (L) = L
		      \]
		      as it preserve the dimension of vector space by injectivity. See that \( \sigma: L \to
		      \sigma(L) \) is also an isomorphism of \( K \)-vector space.
	\end{itemize}
\end{remark}

\begin{eg}
	Have \( \QQ \hookrightarrow \QQ[\sqrt{d}] \) with \( d\in\ZZ \) not a square. Have
	\[
		\begin{aligned}
			\sigma: \QQ[\sqrt{d}]     & \to \QQ[\sqrt{d}]                              \\
			(\sqrt{d})^2              & =d \implies \sigma(\sqrt{d})^2 = \sigma(d) = d \\
			\implies \sigma(\sqrt{d}) & = \pm \sqrt{d}
		\end{aligned}
	\]
	thus there are only two possibilities:
	\begin{itemize}
		\item \( \sigma = \Id \).
		\item \( \sigma(a+b\sqrt{d}) = a - b\sqrt{d} \)
	\end{itemize}
	one may also need to check that this is indeed a ring homomorphism.
\end{eg}

\begin{proposition}
	If \( \qo{L}{K} \) is a finite field extension, then \( G(\qo{L}{K}) \) is a finite group.
\end{proposition}

\begin{proof}
	Let \( a_1,\ldots, a_n \) such that \( L = K(a_1,\ldots,a_n) \). See that \( a_i \) is \textcolor{blue}{algebraic} over
	\( K \) since \( \qo{L}{K} \) is finite. Let
	\[
		f = x^d + c_{d-1} x^{d-1} + \cdots + c_0 \in K[x]
	\]
	be the minimal polynomial of \( a_i \), by definition we shall have
	\[
		a_i^d + c_{d-1} a_i^{d-1} + \cdots + c_0 = 0
	\]
	If then \( \sigma \in G(\qo{L}{K}) \), then
	\[
		\sigma(a_i)^d + c_{d-1} \sigma(a_i)^{d-1} + \cdots + c_0 = \sigma(a_i^d + c_{d-1} a_i^{d-1} +
		\cdots + c_0) = 0
	\]
	\( \implies \sigma(a_i) \) is a root of \( f \), thus there are only finitely many possibilities
	for \( \sigma(a_i) \) to choose. Finitely many possibilities for \( a_i \) and finitely many
	possibilties for each of them to map to with the fact that \( \sigma \) is totally determined by
	\[
		\sigma(a_1),\ldots, \sigma(a_n)
	\]
	then only finitely man such \( \sigma \).
\end{proof}

\begin{remark}
	Suppose \( L = K(a) \) for some \( a \) which is algebraic over \( K \). If \( f \) be the minimal
	polynomial of \( a \implies \deg(f) = [L:K] \). \textcolor{blue}{If \( \sigma \in G(\qo{L}{K}) \implies \sigma(a) \) is a root
		of \( f \implies \# \) distinct roots of \( f \leq \deg(f) \implies \)}
	\[
		\abs{G\bace{\qo{L}{K}}} \leq [L:K]
	\]
	for the equality we need to impose some assumption here. But this in general already give us some
	idea on the cardinality of galois group.
\end{remark}

\begin{eg}
	What is
	\[
		G\bace{\qo{\QQ(\sqrt[3]{2})}{\QQ}}?
	\]
	Consider that \( \QQ(\sqrt[3]{2}) \subseteq \RR \) while the other two roots of \( x^3 - 2 \) do
	not lie in \( \RR \) see that it is only the identity map. In short, the element of the Galois
	group is just permuting the roots within the same minimal polynomial.
\end{eg}

\section{Algebraic Closure}
Recall that a field \( K \) is algebraically closed if \( \forall \; f\in K[x], \deg(f)\geq 1 \),
there exists \( a\in K \), s.t. \( f(a) = 0 \). Which is also equivalent to \( \forall \; f \in K[x], \deg(f) \geq 1 \)
with decomposition
\[
	f = c(x-\lambda_1)\cdots(x-\lambda_n)
\]
with \( c\ne 0 \) and \( \lambda_1,\ldots,\lambda_n \in K \) and \( n = \deg(f) \).

\begin{exercise}
	This is also equivalent to: \( \forall \) field extension \( K \hookrightarrow L \), if \( a \in
	L\smallsetminus K \implies a \) is not algebraic over \( K \).
\end{exercise}

\begin{theorem}
	For every field \( K \), there is an algberaic extension \( K \hookrightarrow \overline{K} \),
	such that \( \overline{K} \) is algebraically closed.
\end{theorem}

\begin{definition}[\textbf{Algebraic Closure}]
	\( \overline{K} \) is called the algebraic closure of \( K \).
\end{definition}
Our goal is to show the uniqueness of \( \overline{K} \).

\begin{proposition}
	\label{prop:extension-alg}
	Given two field extensions:
	% \[
	%   \begin{aligned}
	%     K &\hookrightarrow L \\ 
	%     K &\hookrightarrow M 
	%   \end{aligned}
	% \] 
	\[
		\begin{tikzcd}
			& L \\
			K \arrow[ur, hook] \arrow[dr, hook] & \\
			& M
		\end{tikzcd}
	\]
	such that
	\begin{enumerate}
		\item \( \qo{L}{K} \) is algebraic.
		\item \( M \) is algebraic closed.
	\end{enumerate}
	\( \implies \exists \) morphism of field extensions of \( K \), \( L \to M \).
	\begin{note}
		Such morphism \( L \to K \) is not unique in general.
	\end{note}
\end{proposition}

\begin{proof}
	Define
	\[
		\cP = \{(K',\vp') \; |\; K \hookrightarrow K' \hookrightarrow L \text{ subfield}, \vp': K' \to M
		\text{ morphism of } K \text{-algebra.}\}
	\]
	And we define the order on \( \cP \) by \( (K',\vp')\leq (K'',\vp'') \) if \( K'\subseteq K'' \)
	and \( \vp''|_{K'} = \vp' \). The goal here is to show first that we can apply the Zorn's Lemma on
	\( \cP \).

	\begin{exercise}
		Check that
		\begin{enumerate}
			\item This is indeed an order relation. Meaning that it is reflexive, antisymmetric and
			      transitive.
			\item If \( (K_i,\vp_i)_{i\in I} \) is a family of elements of \( \cP \), s.t. any two of them
			      are comparable, then
			      \[
				      \tilde{K} = \bigcup_{i\in I} K_i \text{ is a subextension of } L
			      \]
			      \( \implies \) If \( \tilde\vp: \tilde{K} \to M \) is defined by \( \tilde{\vp}|_{K_i} = \vp_i \),
			      then
			      \begin{itemize}
				      \item This is well-defined. Meaning the union \( \tilde{K} \) is indeed a field and the
				            map \( \vp \) is well-defined.
				      \item It is a morphism of \( K \)-algebra.
			      \end{itemize}
			\item \( (K_i,\vp_i) \leq (\tilde{K},\tilde\vp) \; \forall \; i \).
		\end{enumerate}
	\end{exercise}
	Now one can apply the \textbf{Zorn's Lemma} \ref{lem:zorn}
	to see that there exists:
	\[
		\begin{tikzcd}
			K \arrow[r,  hook] \arrow[dr] & K_1 \arrow[d, "\vp_1"] \arrow[r,hook] & L \\
			&M
		\end{tikzcd}
	\]
	s.t. \( (K_1,\vp_1) \) is a maximal element of \( \cP \). The goal here is to show that there is
	nothing in between \( K_1 \) and \( L \) otherwise contradicts with the maximality of \(
	(K_1,\vp_1) \).
	\begin{claim}
		\( K_1 = L \) (\( \implies \) QED).
	\end{claim}

	\begin{proof}[\textbf{Proof of Claim}]
		If \( K_1 \ne L \implies \exists \; a \in L \smallsetminus K_1 \), with \( a \) is algebraic
		over \( K \implies a \) is algebraic over \( K_1 \) as \( K_1 \) is the larger field. Let \( f \) be the minimal
		polynomial of \( a \) over \( K_1 \), we have with \( K_1 \subsetneq K_1(a) \)
		\[
			\begin{aligned}
				 & K_1 \hookrightarrow K_1(a) \cong \qo{K_1[x]}{(f)} \\
				 & K_1\xhookrightarrow{\vp_1} M
			\end{aligned}
		\]
		since \( f\in K_1[x] \), \( M \) is algebraically closed, then there exists \( b \in M \), s.t.
		\( f(b) = 0 \). By the universal property of \( K_1[x] \) and of quotient ring, we have
		\[
			\begin{aligned}
				\exists ! \psi: \qo{K_1[x]}{(f)} & \to M     \\
				\overline{a}                     & \mapsto b \\
			\end{aligned}
		\]
		by the universal property that:
		\[
			\begin{tikzcd}
				K_1 \arrow[d, "\vp_1"] \arrow[r, hook, "i_1"] & K_1[x] \arrow[dashed, hook, ld ,"\psi_0"] \arrow[r,
					hook, "i_2"] & \qo{K_1[x]}{(f)} \arrow[lld, hook, dashed, "\psi"] \\
				M & &
			\end{tikzcd}
		\]
		which contradicts to the maximality of \( (K_1,\vp_1) \) \( \lightning \).
	\end{proof}
\end{proof}

\begin{corollary}
	If we have
	% \[
	% 	\begin{aligned}
	% 		K & \hookrightarrow K_1 \\
	% 		K & \hookrightarrow K_2
	% 	\end{aligned}
	% \]
	\[
		\begin{tikzcd}
			K \arrow[r,hook] \arrow[rd,hook] & K_1 \arrow[d, dashed, "\vp"] \\
			& K_2
		\end{tikzcd}
	\]
	are algebraically closures of \( K \), then there exists an \textbf{isomorphism} of \( K \)-extension:
	\[
		\vp: K_1 \cong K_2
	\]
	\begin{note}
		In particular, this yields us the \textcolor{blue}{uniqueness} of algebraic closure up to isomorphism.
	\end{note}
\end{corollary}

\begin{proof}
	We have
	\[
		\begin{cases}
			\qo{K_1}{K} \text{ algebraic} \\
			\qo{K_2}{K} \text{ algebraically closed}
		\end{cases}
		\implies \exists \text{ morphism of } K \text{-extension } \vp: K_1 \hookrightarrow K_2
	\]
	with the implication given by proposition above. Such morphism is clearly injective, and thus we have
	\[
		K\hookrightarrow K_1 \xhookrightarrow{\vp} K_2
	\]
	If there exists \( a \in K_2 \smallsetminus K_1 \), with \( a \) is algebraic over \( K \implies a \)
	is algebraic over \( K_1 \), which contradicts to the fact that \( K_1 \) is algebraically closed
	\( \lightning \).
	Thus \( \vp \) is surjection and being isomorphism.
\end{proof}

\section{Frobenius Homomorphism}
Suppose \( R \) is a ring of characteristic \( p > 0 \) with \( p \) prime integer:
\[
	p \cdot 1_R = 0 \iff \exists \text{ morphism } \qo{\ZZ}{p\ZZ} \hookrightarrow R
\]

\begin{eg}
	Any field \( K \) of characteristic \( p \) or \( K[x_1,\ldots,x_n] \) with \( K \)-field of
	characteristic \( p \) can define the following frobenius homomorphism.
\end{eg}

\begin{proposition}
	\label{prop:frobenius-hom}
	If \( R \) is a commutative ring of \textcolor{blue}{char \( p \)}, \( p \) prime integer. If
	\[
		\begin{aligned}
			F: R & \to R \\
			F(u) & = u^p
		\end{aligned}
	\]
	then \( F \) is a ring homomorphism.
\end{proposition}

\begin{proof}
	Clearly
	\begin{itemize}
		\item \( F(1) = 1 \).
		\item \( F(uv) = F(u)F(v) \).
	\end{itemize}
	And
	\[
		F(a+b) = (a+b)^p = a^p + b^p +\sum_{i=1}^{p-1}\binom{p}{i} a^i b^{p-i} = a^p + b^p = F(a) + F(b)
	\]
	with
	\[
		\binom{p}{i} = \frac{p!}{i!(p-i)!} = \frac{p(p-1)\cdots(p-i+1)}{i!} \equiv 0 \; (\bmod{p}) \quad
		\forall \; 1
		\leq i \leq p-1
	\]
	thus vanishes.
\end{proof}
